
\documentclass[12pt,a4paper]{article}
\usepackage[utf8]{inputenc}
\usepackage[T1]{fontenc}
\usepackage[brazil]{babel}
\usepackage{setspace}
\usepackage{geometry}
\geometry{a4paper, margin=2.5cm}
\setstretch{1.3}

\title{Referencial Teórico e Metodologia do Sistema de Treinamento com Mutação e Regra de Classificação}
\author{Projeto QIP}
\date{\today}

\begin{document}
\maketitle

\section*{Referencial Teórico}

O objetivo deste trabalho é explicar, de maneira acessível até mesmo para pessoas sem formação técnica,
como funciona um sistema de classificação baseado em aprendizado heurístico.
De forma simplificada, o que se deseja é ensinar o computador a ``pesar'' diferentes características
(dados de entrada) para chegar a uma probabilidade de cada possível resultado (classes).

\subsection*{Classificação de Dados}
Imagine que temos várias observações (como respostas de questionários ou medições clínicas) e queremos prever
se cada pessoa pertence a uma ou mais categorias (como diferentes diagnósticos). O sistema analisa as colunas
de dados e tenta aprender quais delas são mais importantes para cada categoria.

\subsection*{Probabilidades e Softmax}
O computador transforma a soma ponderada de cada característica em uma probabilidade.
Para isso, utiliza-se uma função chamada \emph{softmax}, que garante que todas as probabilidades fiquem
entre 0 e 1 e somem exatamente 1. Assim, podemos dizer, por exemplo, que para uma linha de dados a probabilidade
é 70\% de pertencer à Classe A, 20\% à Classe B e 10\% à Classe C.

\subsection*{Treino, Validação e Desconhecidos}
Os dados foram separados em três grupos:
\begin{itemize}
    \item \textbf{Treino:} usado para ajustar os pesos do modelo;
    \item \textbf{Validação:} usado para verificar se o ajuste está realmente funcionando em dados diferentes dos usados no treino;
    \item \textbf{Desconhecidos:} linhas em que o alvo era ``não'' (ou vazio), ou seja, casos em que não sabemos a classe real. Eles não participam do treino nem da validação, mas recebem uma classificação final.
\end{itemize}

\subsection*{Classe ``Sem Transtorno''}
Existe uma categoria especial chamada \textbf{Sem Transtorno}.
Ela não é aprendida diretamente no treino, mas surge de uma \emph{regra} que observa quando o modelo está inseguro
(em casos em que a maior probabilidade é baixa e a diferença para a segunda maior também é pequena). Nestes casos,
o sistema aloca parte da probabilidade para a classe ``Sem Transtorno''.

\subsection*{Mutação e Otimização}
Para melhorar a busca por bons pesos, além do ajuste tradicional (gradiente), o sistema aplica \textbf{mutações aleatórias}
em cinco colunas de cada vez. Isso ajuda a escapar de soluções ruins e explorar alternativas.
Outras técnicas também são usadas para acelerar e estabilizar o aprendizado, como:
\begin{itemize}
    \item \textbf{Lookahead:} suaviza oscilações juntando soluções rápidas com uma versão mais estável;
    \item \textbf{EMA (média exponencial):} mantém uma média dos melhores pesos ao longo do tempo;
    \item \textbf{Pesos por classe:} dá mais importância a classes raras para não deixá-las esquecidas;
    \item \textbf{Annealing:} começa mais flexível e vai ficando mais rígido com o tempo;
    \item \textbf{Simulated Annealing:} permite aceitar pequenas piores soluções no início para não ficar preso em um ponto ruim;
    \item \textbf{Dropout:} desliga temporariamente algumas colunas no treino, evitando dependência excessiva;
    \item \textbf{Barzilai--Borwein:} ajusta automaticamente o tamanho do passo do aprendizado.
\end{itemize}

\section*{Metodologia}

\subsection*{Etapas do Processo}
O processo completo seguiu as etapas a seguir:

\begin{enumerate}
    \item \textbf{Preparação dos Dados:} leitura da planilha Excel, separando colunas de características (entradas) e a coluna Alvo (saídas desejadas).
    Linhas com alvo ``não'' ou vazio foram marcadas como \emph{desconhecidas}.
    \item \textbf{Definição das Classes:} foram consideradas as classes originais (como diferentes transtornos) e uma classe adicional ``Sem Transtorno'', que é ativada por regra e não por treino direto.
    \item \textbf{Divisão dos Conjuntos:} os casos conhecidos foram divididos em treino e validação, garantindo que todas as classes estivessem representadas em ambos os grupos.
    \item \textbf{Treinamento com Mutação:} o modelo foi ajustado por descida de gradiente proximal (técnica que combina aproximação suave e regularização), aliado a mutações em 5 colunas de cada vez. Foram aplicados os métodos auxiliares (Lookahead, EMA, etc.) para melhorar a convergência.
    \item \textbf{Interrupção Manual:} o treinamento ocorre indefinidamente e só é interrompido quando o usuário pressiona a tecla ``q''. Nesse momento, são registrados os melhores parâmetros encontrados até então.
    \item \textbf{Validação:} após o fim do treino, o modelo é avaliado no conjunto de validação, gerando métricas de desempenho (como macro top-3, F1 e perda de Hamming).
    \item \textbf{Etapa Final:} são calculadas as probabilidades para todos os casos (incluindo os desconhecidos). A regra de ``Sem Transtorno'' é aprendida usando todas as linhas, definindo os parâmetros T1, T2 e $\gamma$. É gerada a aba ``Resultado\_Final'' na planilha com as probabilidades e o ranking top-3 para cada linha.
\end{enumerate}

\subsection*{Saídas Produzidas}
O sistema grava várias abas no Excel, entre elas:
\begin{itemize}
    \item \textbf{Pontuação\_Tunada:} pesos ajustados por classe e característica;
    \item \textbf{Resultado\_Heuristica\_Tunada:} probabilidades no treino;
    \item \textbf{Metricas\_Heuristica\_Tunada:} métricas no treino;
    \item \textbf{Resultado\_Validacao:} probabilidades no conjunto de validação;
    \item \textbf{Metricas\_Validacao:} métricas no conjunto de validação;
    \item \textbf{Split\_Info:} resumo de como foi feita a divisão treino/validação;
    \item \textbf{Regras\_Normal:} parâmetros T1, T2, $\gamma$ aprendidos;
    \item \textbf{Resultado\_Final:} probabilidades e top-3 para todas as linhas (incluindo as desconhecidas).
\end{itemize}

\end{document}
